\documentclass[11pt]{article}

\usepackage[T1]{fontenc}
\usepackage[utf8]{inputenc}
\usepackage{lmodern}
\usepackage[margin=1in]{geometry}
\usepackage{microtype}
\usepackage{amsmath,amssymb,mathtools,bm}
\usepackage{booktabs,tabularx,array}
\usepackage{enumitem}
\usepackage{xcolor}
\usepackage{listings}
\usepackage{hyperref}

\definecolor{codebg}{RGB}{247,247,247}
\definecolor{codeframe}{RGB}{205,205,205}
\definecolor{codecomment}{RGB}{0,105,75}
\definecolor{codekeyword}{RGB}{25,65,145}

\lstdefinestyle{base}{
  basicstyle=\ttfamily\footnotesize,
  columns=fullflexible,
  keepspaces=true,
  showstringspaces=false,
  breaklines=true,
  breakatwhitespace=false,
  frame=single,
  rulecolor=\color{codeframe},
  backgroundcolor=\color{codebg},
  numbers=left,
  numberstyle=\tiny\color{gray},
  numbersep=7pt,
  xleftmargin=1.5em,
  framexleftmargin=1em,
  aboveskip=0.8em,
  belowskip=0.8em
}

\lstdefinestyle{cpp}{
  style=base,
  language=C++,
  keywordstyle=\color{codekeyword}\bfseries,
  commentstyle=\color{codecomment},
  stringstyle=\color{purple!65!black}
}

\lstdefinestyle{data}{
  style=base,
  language={}
}

\hypersetup{
  colorlinks=true,
  linkcolor=blue!55!black,
  urlcolor=blue!65!black,
  citecolor=blue!55!black,
  pdftitle={BCC Binary-Junction Node Splitting in ParaDiS}
}

\setlist{nosep,leftmargin=*}
\setcounter{tocdepth}{2}
\allowdisplaybreaks
\newcommand{\code}[1]{\texttt{\detokenize{#1}}}
\newcommand{\vect}[1]{\bm{#1}}

\title{BCC Binary-Junction Node Splitting in ParaDiS\\
       \large File-by-File Code and Design Guide}
\author{ParaDiS-Extended Development Notes}
\date{August 2026}

\begin{document}

\maketitle

\begin{abstract}
This note documents the implementation of three-arm BCC binary-junction
splitting in the serial and parallel ParaDiS topology paths.  Unlike a
high-level feature summary, it shows the code added or changed in every
implementation and regression file and explains the purpose of each change.
The port follows the ExaDiS/OpenDiS implementation pinned at commit
\code{351607fa0f04aa34d29ee253438edde994473fd6}, with ParaDiS-specific
adaptations for node constraints, remeshing, and serialized cross-domain
topology operations.
\end{abstract}

\tableofcontents
\newpage

\section{Scope, lineage, and reading convention}

The functional implementation is the difference from the pre-feature base
\code{efd4069bfb9540a2ed76ccbe9892ed2c441597f7} through
\code{3cb6c8bd92d7407a5b80b595d116574dc6b86f7f}.  The first implementation
commit is \code{74f18516eb24587c65d123e77a9713c71f38d2f0}; the second completes MPI
constraint propagation, remesh protection, and the regression fixture.  The
exact change set can be inspected with:

\begin{lstlisting}[style=data]
git diff efd4069bfb9540a2ed76ccbe9892ed2c441597f7 \
         3cb6c8bd92d7407a5b80b595d116574dc6b86f7f -- \
         include src tests
\end{lstlisting}

The pinned reference is available at:

\begin{center}
\url{https://github.com/llnl/exadis/tree/351607fa0f04aa34d29ee253438edde994473fd6/src/topology_types}
\end{center}

Every changed file has a subsection below.  Listings reproduce the final
code for the introduced behavior, with only enough surrounding context to
make each fragment understandable.  Large functions are divided into focused
fragments rather than copied wholesale.  Therefore the listings are a design
and change guide, while the repository source remains authoritative.  Long
source lines are occasionally reindented to fit the page without changing
their meaning.  Comments beginning with ``Existing'' are editorial elisions,
not new comments in the C++ source.

The 17 changed implementation and test files fall into five groups:

\begin{itemize}
  \item public definitions: four headers under \path{include};
  \item parameter and initialization plumbing: four source files;
  \item serial and parallel topology evaluation: four source files;
  \item remesh and MPI propagation: three source files; and
  \item regression inputs: one control file and one nodal-data file.
\end{itemize}

\section{Public definitions and interfaces}

\subsection{\texttt{include/Node.h}}

\paragraph{New lines, with existing enumeration entries shown for context.}

\begin{lstlisting}[style=cpp]
#define CORNER_NODE         0x00000010   /* 16: Preserve a physical corner */
                                         /*     during mesh coarsening.     */
\end{lstlisting}

\paragraph{Meaning.}
The new constraint bit marks the degree-two node created by a successful
binary-junction split.  It is a topology state, not a temporary flag: remesh
uses it to distinguish a physical kink from an ordinary degree-two node that
may be coarsened away.  Decimal value 16 is deliberately disjoint from the
existing pinning and surface bits.

\subsection{\texttt{include/Param.h}}

\paragraph{Code added.}

\begin{lstlisting}[style=cpp]
int split3node;                /* Toggle BCC binary-junction */
                               /* three-arm node splitting. */
\end{lstlisting}

\paragraph{Meaning.}
This member carries the runtime control-file switch through \code{Param_t}.
It lets a user enable or disable degree-three BCC splitting without changing
the established treatment of nodes with four or more arms.

\subsection{\texttt{include/Topology.h}}

\paragraph{Code added.}

\begin{lstlisting}[style=cpp]
int BCC_binary_junction_node(Home_t *home, Node_t *node,
        real8 tjunc[3], int *planarJunc);
\end{lstlisting}

\paragraph{Meaning.}
The declaration exposes one shared classifier to the serial topology path and
the parallel candidate/force stages.  The return value is the junction-arm
index or $-1$.  The output vector \code{tjunc} is the unit junction-line
direction, and \code{planarJunc} reports whether the parent-plane geometry is
planar and therefore excluded from this split mechanism.

\subsection{\texttt{include/OpList.h}}

\paragraph{Code added.}

\begin{lstlisting}[style=cpp]
typedef enum {
        REM_OP_CHANGE_CONN,
        REM_OP_INSERT_ARM,
        REM_OP_REMOVE_NODE,
        REM_OP_CHANGE_BURG,
        REM_OP_SPLIT_NODE,
        REM_OP_RESET_COORD,
        REM_OP_RESET_SEG_FORCE1,
        REM_OP_RESET_SEG_FORCE2,
        REM_OP_MARK_FORCE_OBSOLETE,
        REM_OP_RESET_PLANE,
        REM_OP_RESET_CONSTRAINT
} OpType_t;

typedef struct {
    int   opType;
    Tag_t tag1;
    int   constraint;
} RemOpResetConstraint_t;

void AddOpResetConstraint(Home_t *home, Tag_t *tag1, int constraint);
\end{lstlisting}

\paragraph{Meaning.}
ParaDiS communicates topology changes as packed operation records.  The new
record transmits a node tag plus the complete integer constraint bitmask.  Its
operation identifier is appended to the enumeration, preserving the numeric
values of all existing operation types.  All MPI ranks must still run the
same rebuilt executable because the byte-stream protocol now includes this
record.

\section{Parameter and initialization plumbing}

\subsection{\texttt{src/Initialize.cc}}

\paragraph{Code changed.}

\begin{lstlisting}[style=cpp]
allValidConstraints = PINNED_X | PINNED_Y | PINNED_Z |
                      SURFACE_NODE | CORNER_NODE |
                      MULTI_JUNCTION_NODE |
                      LOOP_NODE | CROSS_NODE;
\end{lstlisting}

\paragraph{Meaning.}
Restart and input validation now accepts constraint bit 16.  Without this
addition, a restart written after a binary-junction split would be rejected as
containing an unknown node constraint.

\subsection{\texttt{src/InputSanity.cc}}

\paragraph{Code added.}

\begin{lstlisting}[style=cpp]
param->split3node = (param->split3node != 0);
\end{lstlisting}

\paragraph{Meaning.}
Any nonzero input is normalized to one.  Downstream candidate logic can treat
\code{split3node} as a strict Boolean rather than preserving arbitrary
integer values from a control file.

\subsection{\texttt{src/ParadisInit.cc}}

\paragraph{Code added.}

\begin{lstlisting}[style=cpp]
printf("|   split BCC three-arm nodes    : %s\n",
       p->split3node ? "yes" : "no");
\end{lstlisting}

\paragraph{Meaning.}
The startup summary makes the effective state visible before stepping begins.
This is especially useful when comparing otherwise identical serial-selector
and parallel-selector regression runs.

\subsection{\texttt{src/Param.cc}}

\paragraph{Code added or changed.}

\begin{lstlisting}[style=cpp]
BindVar(CPList, "split3node", &param->split3node,
        V_INT, 1, VFLAG_NULL);
param->split3node = 1;

printf("   splitMultiNodeAlpha                  : %+lf\n",
       p->splitMultiNodeAlpha);
printf("   split3node                           : %d\n\n",
       p->split3node);
\end{lstlisting}

\paragraph{Meaning.}
\code{CtrlParamInit()} registers the control keyword and enables it by
default.  \code{Print_Param()} includes the value in the detailed parameter
dump.  Existing control files therefore gain the feature automatically, while
\code{split3node = 0} restores the former degree-four minimum.

\section{Shared classifier and serial topology path}

\subsection{\texttt{src/Topology.cc}}

This file contains the shared BCC classifier and the serial implementation in
\code{SplitMultiNodes()}.

\subsubsection{Removing stale corner state}

\paragraph{Code added.}

\begin{lstlisting}[style=cpp]
if (HAS_ANY_OF_CONSTRAINTS(localNode->constraint, CORNER_NODE) &&
    (localNode->numNbrs != 2)) {
    REMOVE_CONSTRAINTS(localNode->constraint, CORNER_NODE);
}

if (HAS_ANY_OF_CONSTRAINTS(ghostNode->constraint, CORNER_NODE) &&
    (ghostNode->numNbrs != 2)) {
    REMOVE_CONSTRAINTS(ghostNode->constraint, CORNER_NODE);
}
\end{lstlisting}

\paragraph{Meaning.}
\code{InitTopologyExemptions()} enforces the invariant that a corner marker is
meaningful only on a degree-two node.  If a later collision or topology event
changes the degree, native and ghost copies lose stale corner state during the
next \code{InitTopologyExemptions()} pass.

\subsubsection{Recognizing a BCC binary junction}

\paragraph{Code added or moved, with both adjustment calls shown for
context.}

\begin{lstlisting}[style=cpp]
int BCC_binary_junction_node(Home_t *home, Node_t *node,
                             real8 tjunc[3], int *planarJunc)
{
        int     armID;
        int     binaryJunc;
        int     numArmEdge, numArmGlide, numArmJunc;
        real8   ndir[3][3];
        real8   eps, angleTol;
        Param_t *param;

        binaryJunc = -1;
        numArmEdge = 0;
        numArmGlide = 0;
        numArmJunc = 0;

        VECTOR_ZERO(tjunc);
        *planarJunc = 0;

        if ((node == (Node_t *)NULL) || (node->numNbrs != 3)) {
            return(-1);
        }

        param = home->param;
        eps = 1.0e-12;
        angleTol = 5.0 * M_PI / 180.0;

        for (armID = 0; armID < node->numNbrs; armID++) {
            real8  bMag, lineMag2;
            real8  burg[3], lineDir[3], normal[3];
            Node_t *nbrNode;

            nbrNode = GetNeighborNode(home, node, armID);

            if (nbrNode == (Node_t *)NULL) {
                continue;
            }

            lineDir[X] = nbrNode->x - node->x;
            lineDir[Y] = nbrNode->y - node->y;
            lineDir[Z] = nbrNode->z - node->z;

            ZImage(param, &lineDir[X], &lineDir[Y], &lineDir[Z]);
            lineMag2 = DotProduct(lineDir, lineDir);

            if (lineMag2 < eps) {
                continue;
            }

            NormalizeVec(lineDir);

            burg[X] = node->burgX[armID];
            burg[Y] = node->burgY[armID];
            burg[Z] = node->burgZ[armID];
            bMag = Normal(burg);

            if (bMag > (1.0 + eps)) {
                VECTOR_COPY(tjunc, lineDir);
                binaryJunc = armID;
                numArmJunc++;
            } else {
                CrossVector(burg, lineDir, normal);

                if (fabs(DotProduct(normal, normal)) > 1.0e-2) {
                    VECTOR_COPY(ndir[numArmEdge], normal);
                    numArmEdge++;
                }

                numArmGlide++;
            }
        }

        if ((numArmJunc == 1) && (numArmGlide == 2)) {
            if (numArmEdge == 2) {
                real8 collinearTol, tdir[3];

                collinearTol = 1.0 - cos(angleTol);
                CrossVector(ndir[0], ndir[1], tdir);
                NormalizeVec(tdir);

                *planarJunc =
                    (fabs(fabs(DotProduct(tdir, tjunc)) - 1.0) <
                     collinearTol);
            } else if (numArmEdge == 1) {
                *planarJunc =
                    (fabs(DotProduct(ndir[0], tjunc)) < sin(angleTol));
            } else {
                *planarJunc = 1;
            }
        } else {
            binaryJunc = -1;
        }

        return(binaryJunc);
}
\end{lstlisting}

\paragraph{Meaning.}
The classifier requires exactly one arm with
$\lVert\vect b\rVert>1+10^{-12}$ and two remaining non-junction arms.  Each
line direction is minimum-imaged and normalized.  For the two non-junction
arms it constructs $\vect n=\vect b\times\vect t$ and treats
$\lVert\vect n\rVert^2>10^{-2}$ as edge-bearing.  A five-degree tolerance
then tests whether the junction direction follows the intersection of two
parent planes, lies in the sole edge parent plane, or has two screw parents.
Those cases are planar and are not split.  The magnitude test calls the two
remaining arms ``glide'' arms in the implementation; it does not independently
validate a glissile Burgers-vector family.

\subsubsection{Admitting the serial degree-three candidate}

\paragraph{Code added or changed.}

\begin{lstlisting}[style=cpp]
int binaryJunc = -1;

if (HAS_ANY_OF_CONSTRAINTS(node->constraint, CORNER_NODE) &&
    (node->numNbrs != 2)) {
    REMOVE_CONSTRAINTS(node->constraint, CORNER_NODE);
}

if (node->numNbrs < 3) continue;

if (node->numNbrs == 3) {
    int   planarJunc;
    real8 tjunc[3];

    if ((!param->split3node) ||
        (param->materialType != MAT_TYPE_BCC)) {
        continue;
    }

    binaryJunc = BCC_binary_junction_node(home, node, tjunc,
                                           &planarJunc);

    if ((binaryJunc < 0) || planarJunc) continue;
}

if (HAS_ANY_OF_CONSTRAINTS(node->constraint, SURFACE_NODE)) {
    continue;
}

if (HAS_ANY_OF_CONSTRAINTS(node->constraint, PINNED_NODE)) {
    continue;
}
\end{lstlisting}

\paragraph{Meaning.}
The serial candidate floor changes from degree four to degree three, but the
new degree-three path is narrow: it requires the switch, BCC material, a
recognized non-planar binary junction, and an unpinned/non-surface node.  The
existing short-segment and cycle-frequency checks still apply later in the
function.

\subsubsection{Selecting only the two non-junction arms}

\paragraph{Code changed.}

\begin{lstlisting}[style=cpp]
armList  = (int *)calloc(1, (nbrs - 1) * sizeof(int));
armList2 = (int *)calloc(1, (nbrs - 1) * sizeof(int));

for (j = 0; j < numSets; j++) {
    int tmpRepositionNode1, tmpRepositionNode2;

    if ((nbrs == 3) && (armSets[j][binaryJunc] != 0)) {
        continue;
    }

    /* Existing trial-split evaluation follows. */
}
\end{lstlisting}

\paragraph{Meaning.}
Two selected arms must fit in the work arrays, hence
\code{nbrs - 1} instead of \code{nbrs - 2}.  The pre-existing split-count
table already contains three combinations for a degree-three node.  The new
filter discards the two combinations containing the junction arm, leaving
only the trial that moves both non-junction arms.

\subsubsection{Degree-two mobility boundary condition}

\paragraph{Code added.}

\begin{lstlisting}[style=cpp]
AdjustJuncNodeForceAndVel(home, splitNode1, &mobArgs1);
AdjustJuncNodeForceAndVel(home, splitNode2, &mobArgs2);

if (nbrs == 3) {
    Node_t *twoNode =
        (splitNode1->numNbrs == 2) ? splitNode1 :
        ((splitNode2->numNbrs == 2) ? splitNode2 :
         (Node_t *)NULL);

    if (twoNode != (Node_t *)NULL) {
        twoNode->vX = 0.0;
        twoNode->vY = 0.0;
        twoNode->vZ = 0.0;
    }
}
\end{lstlisting}

The same zeroing block is applied again after the repositioned trial has its
mobility recomputed.

\paragraph{Meaning.}
Both trial nodes receive the established junction force/velocity adjustment.
For this binary split, the resulting degree-two corner is stationary during
the energetic test, matching the reference boundary condition and preventing
it from contributing artificial power.

\subsubsection{Guarding direction and accepting power}

\paragraph{Code added or changed.}

\begin{lstlisting}[style=cpp]
if (((nbrs != 3) || !param->enforceGlidePlanes) &&
    (vd1 > eps) && (vd2 > eps) &&
    (fabs(vel1Dotvel2/vd1Sqrt/vd2Sqrt+1.0) < 0.01)) {
    /* Existing equal-and-opposite velocity handling. */
}

if (((powerTest - powerMax) > eps) &&
    ((nbrs != 3) || (powerTest > 1.0))) {
    vel1[X] = splitNode1->vX;
    vel1[Y] = splitNode1->vY;
    vel1[Z] = splitNode1->vZ;

    vel2[X] = splitNode2->vX;
    vel2[Y] = splitNode2->vY;
    vel2[Z] = splitNode2->vZ;

    splitPos1[X] = splitNode1->x;
    splitPos1[Y] = splitNode1->y;
    splitPos1[Z] = splitNode1->z;

    splitPos2[X] = splitNode2->x;
    splitPos2[Y] = splitNode2->y;
    splitPos2[Z] = splitNode2->z;
}
\end{lstlisting}

\paragraph{Meaning.}
The norm tests prevent division by zero when testing nearly opposite
velocities.  Degree-three handling also respects enforced glide planes.  A
binary-junction trial must improve on the current best power and exceed the
absolute threshold \code{powerTest > 1.0}.  Saving the winning velocities and
trial positions already existed; the new degree-three behavior is the stricter
acceptance condition and the later reuse of those saved positions.  The serial
path already used
$d_{\mathrm{split}}=2r_{\mathrm{ann}}+10^{-12}$ before this feature; that line
is supporting context, not a newly introduced serial change.

\subsubsection{Using the evaluated position and preserving the corner}

\paragraph{Code added.}

\begin{lstlisting}[style=cpp]
if (nbrs == 3) {
    VECTOR_COPY(pos1, splitPos1);
    VECTOR_COPY(pos2, splitPos2);
}

/* SplitNode() executes the accepted final split here. */

if (nbrs == 3) {
    if (splitNode1->numNbrs == 2) {
        ADD_CONSTRAINTS(splitNode1->constraint, CORNER_NODE);
        if (globalOp) {
            AddOpResetConstraint(home, &splitNode1->myTag,
                                 splitNode1->constraint);
        }
    }
    if (splitNode2->numNbrs == 2) {
        ADD_CONSTRAINTS(splitNode2->constraint, CORNER_NODE);
        if (globalOp) {
            AddOpResetConstraint(home, &splitNode2->myTag,
                                 splitNode2->constraint);
        }
    }
}
\end{lstlisting}

\paragraph{Meaning.}
The final degree-three split reuses the positions that actually passed the
trial mobility/power test instead of entering the general
\code{CalcJunctionLen()} path.  After \code{SplitNode()} succeeds, whichever
output node has degree two receives \code{CORNER_NODE}.  A global operation
then appends the full constraint reset after the create/connectivity records.

\section{Parallel candidate and force evaluation}

\subsection{\texttt{src/SMNSupport.cc}}

\paragraph{Code added: shared parallel eligibility predicate.}

\begin{lstlisting}[style=cpp]
static int IsSplittableMultiNode(Home_t *home, Node_t *node)
{
        int   binaryJunc, planarJunc;
        real8 tjunc[3];

        if (node->numNbrs < 3) {
            return(0);
        }

        if (node->numNbrs > 3) {
            return(1);
        }

        if (HAS_ANY_OF_CONSTRAINTS(node->constraint,
                                   SURFACE_NODE | PINNED_NODE)) {
            return(0);
        }

        if ((!home->param->split3node) ||
            (home->param->materialType != MAT_TYPE_BCC)) {
            return(0);
        }

        binaryJunc = BCC_binary_junction_node(home, node, tjunc,
                                               &planarJunc);

        return((binaryJunc >= 0) && !planarJunc);
}
\end{lstlisting}

\paragraph{Code changed: degree-three enumeration and lists.}

\begin{lstlisting}[style=cpp]
maxSets = POSSIBLE_SPLITS[numNbrs];
maxSplitCnt = numNbrs >> 1;

if (numNbrs == 3) {
    maxSplitCnt = 2;
}

/* BuildRecvDomList(): */
if (lmnInfo[i].multiNode->numNbrs < 3) {
    continue;
}

/* Applied in both local and ghost list builders: */
if (HAS_ANY_OF_CONSTRAINTS(node->constraint, CORNER_NODE) &&
    (node->numNbrs != 2)) {
    REMOVE_CONSTRAINTS(node->constraint, CORNER_NODE);
}

if (!IsSplittableMultiNode(home, node)) {
    continue;
}
\end{lstlisting}

\paragraph{Meaning.}
Native and ghost lists now agree on the exact degree-three eligibility rules.
\code{SMNGetArmSets()} is allowed to select two arms from a three-arm node;
the existing \code{POSSIBLE_SPLITS[3] == 3} table entry supplies the three
combinations.  Receive-domain construction also admits degree three, and both
list builders clear stale corner state before applying the predicate.  The
unchanged \path{src/SplitMultiNodesParallel.cc} then orchestrates these lists,
partial-force communication, and final owning-domain evaluation.

\subsection{\texttt{src/SMNEvalSplitForces.cc}}

This stage evaluates local force contributions for native and ghost
candidates before point-to-point force exchange and accumulation on the
owning domain.

\paragraph{Code changed: distance, classification, and arm selection.}

\begin{lstlisting}[style=cpp]
eps = 1.0e-12;
splitDist = param->rann * 2.0 + eps;

int binaryJunc = -1;

if (numNbrs == 3) {
    int   planarJunc;
    real8 tjunc[3];

    binaryJunc = BCC_binary_junction_node(home, node, tjunc,
                                           &planarJunc);
    if (planarJunc) binaryJunc = -1;
}

armList = (int *)calloc(1, (numNbrs - 1) * sizeof(int));

for (setID = 0; setID < totalSets; setID++) {
    int tmpRepositionNode1 = 0, tmpRepositionNode2 = 0;

    if ((numNbrs == 3) &&
        ((binaryJunc < 0) ||
         (armSets[setID][binaryJunc] != 0))) {
        continue;
    }

    /* Existing partial-force trial follows. */
}
\end{lstlisting}

\paragraph{Code changed: mobility status and degree-two velocity.}

\begin{lstlisting}[style=cpp]
mobStatus  = EvaluateMobility(home, splitNode1, &mobArgs1);
mobStatus |= EvaluateMobility(home, splitNode2, &mobArgs2);

AdjustJuncNodeForceAndVel(home, splitNode1, &mobArgs1);
AdjustJuncNodeForceAndVel(home, splitNode2, &mobArgs2);

if (numNbrs == 3) {
    Node_t *twoNode =
        (splitNode1->numNbrs == 2) ? splitNode1 :
        ((splitNode2->numNbrs == 2) ? splitNode2 :
         (Node_t *)NULL);

    if (twoNode != (Node_t *)NULL) {
        twoNode->vX = 0.0;
        twoNode->vY = 0.0;
        twoNode->vZ = 0.0;
    }
}

if (((numNbrs != 3) || !param->enforceGlidePlanes) &&
    (vd1 > eps) && (vd2 > eps) &&
    (fabs(vel1Dotvel2/vd1Sqrt/vd2Sqrt+1.0) < 0.01)) {
    /* Existing equal-and-opposite velocity handling. */
}
\end{lstlisting}

\paragraph{Meaning.}
The positive epsilon places the trial connector just beyond the immediate
corner-remesh threshold.  Reclassification protects against stale list data,
and only the arm set excluding the junction arm is evaluated.  Work storage
fits two selected arms and reposition flags start deterministically at zero.
Using \code{|=} retains an error from either mobility call.  Both trial nodes
are adjusted, the degree-two node is stationary, and velocity normalization
cannot divide by a zero norm.

\subsection{\texttt{src/SMNEvalMultiNodeSplits.cc}}

This owning-domain stage receives the complete forces, repeats the mobility
and power test, and executes the accepted topology change.

\paragraph{Code changed: candidate parity with the force stage.}

\begin{lstlisting}[style=cpp]
eps       = 1.0e-12;
vNoise    = param->splitMultiNodeAlpha * param->rTol / param->deltaTT;
splitDist = param->rann * 2.0 + eps;

int binaryJunc = -1;

if (numNbrs == 3) {
    int   planarJunc;
    real8 tjunc[3];

    binaryJunc = BCC_binary_junction_node(home, node, tjunc,
                                           &planarJunc);
    if (planarJunc) binaryJunc = -1;
}

armList  = (int *)calloc(1, (numNbrs - 1) * sizeof(int));
armList2 = (int *)calloc(1, (numNbrs - 1) * sizeof(int));

for (setID = 0; setID < totalSets; setID++) {
    int tmpRepositionNode1 = 0, tmpRepositionNode2 = 0;

    if ((numNbrs == 3) &&
        ((binaryJunc < 0) ||
         (armSets[setID][binaryJunc] != 0))) {
        continue;
    }

    /* Existing final trial evaluation follows. */
}
\end{lstlisting}

\paragraph{Code added: degree-two boundary condition.}

\begin{lstlisting}[style=cpp]
AdjustJuncNodeForceAndVel(home, splitNode1, &mobArgs1);
AdjustJuncNodeForceAndVel(home, splitNode2, &mobArgs2);

if (numNbrs == 3) {
    Node_t *twoNode =
        (splitNode1->numNbrs == 2) ? splitNode1 :
        ((splitNode2->numNbrs == 2) ? splitNode2 :
         (Node_t *)NULL);

    if (twoNode != (Node_t *)NULL) {
        twoNode->vX = 0.0;
        twoNode->vY = 0.0;
        twoNode->vZ = 0.0;
    }
}
\end{lstlisting}

This block occurs after both the coincident trial mobility and the
repositioned trial mobility.

\paragraph{Code changed: direction guard, power, final positions, and
constraint.}

\begin{lstlisting}[style=cpp]
if (((numNbrs != 3) || !param->enforceGlidePlanes) &&
    (vd1 > eps) && (vd2 > eps) &&
    (fabs(vel1Dotvel2/vd1Sqrt/vd2Sqrt+1.0) < 0.01)) {
    /* Existing equal-and-opposite velocity handling. */
}

if (((powerTest - powerMax) > eps) &&
    ((numNbrs != 3) || (powerTest > 1.0))) {
    /* Save winning velocities and splitPos1/splitPos2. */
}

if (numNbrs == 3) {
    VECTOR_COPY(pos1, splitPos1);
    VECTOR_COPY(pos2, splitPos2);
}

/* SplitNode() executes the accepted final split here. */

if (numNbrs == 3) {
    if (splitNode1->numNbrs == 2) {
        ADD_CONSTRAINTS(splitNode1->constraint, CORNER_NODE);
        if (globalOp) {
            AddOpResetConstraint(home, &splitNode1->myTag,
                                 splitNode1->constraint);
        }
    }
    if (splitNode2->numNbrs == 2) {
        ADD_CONSTRAINTS(splitNode2->constraint, CORNER_NODE);
        if (globalOp) {
            AddOpResetConstraint(home, &splitNode2->myTag,
                                 splitNode2->constraint);
        }
    }
}
\end{lstlisting}

\paragraph{Meaning.}
The owning-domain decision mirrors the serial decision: the degree-two trial
is stationary, power must both improve and exceed one, and the final split
uses the tested positions.  As in the serial function, storing the winning
positions was pre-existing; the new code applies the absolute degree-three
threshold and explicitly copies those positions into the final degree-three
split.  Successful corner marking is immediately followed by a remote
constraint operation, so serial and parallel topology selectors produce the
same persistent topology state.

\section{Remesh preservation and MPI propagation}

\subsection{\texttt{src/Remesh.cc}}

Both remesh rules reach the shared \code{MeshCoarsen()} implementation, so the
corner rule is implemented once in its candidate and execution phases.

\paragraph{Code added: candidate-phase guard.}

\begin{lstlisting}[style=cpp]
if ((HAS_ANY_OF_CONSTRAINTS(node->constraint, CORNER_NODE) ||
     HAS_ANY_OF_CONSTRAINTS(nbr1->constraint, CORNER_NODE)) &&
    (r1 > (2.0 * param->rann)) &&
    (HAS_ANY_OF_CONSTRAINTS(node->constraint, CORNER_NODE) ||
     HAS_ANY_OF_CONSTRAINTS(nbr2->constraint, CORNER_NODE)) &&
    (r2 > (2.0 * param->rann))) {
    continue;
}
\end{lstlisting}

\paragraph{Code added: independent execution-phase gates.}

\begin{lstlisting}[style=cpp]
mergeNbr1Allowed = 1;
mergeNbr2Allowed = 1;

if (HAS_ANY_OF_CONSTRAINTS(node->constraint, CORNER_NODE) ||
    HAS_ANY_OF_CONSTRAINTS(nbr1->constraint, CORNER_NODE)) {
    real8 dx1, dy1, dz1;
    real8 cornerMergeDist2;

    dx1 = nbr1->x - node->x;
    dy1 = nbr1->y - node->y;
    dz1 = nbr1->z - node->z;

    ZImage(param, &dx1, &dy1, &dz1);
    cornerMergeDist2 = 4.0 * param->rann * param->rann;

    mergeNbr1Allowed = ((dx1*dx1 + dy1*dy1 + dz1*dz1) <=
                        cornerMergeDist2);
}

if (HAS_ANY_OF_CONSTRAINTS(node->constraint, CORNER_NODE) ||
    HAS_ANY_OF_CONSTRAINTS(nbr2->constraint, CORNER_NODE)) {
    real8 dx2, dy2, dz2;
    real8 cornerMergeDist2;

    dx2 = nbr2->x - node->x;
    dy2 = nbr2->y - node->y;
    dz2 = nbr2->z - node->z;

    ZImage(param, &dx2, &dy2, &dz2);
    cornerMergeDist2 = 4.0 * param->rann * param->rann;

    mergeNbr2Allowed = ((dx2*dx2 + dy2*dy2 + dz2*dz2) <=
                        cornerMergeDist2);
}

if ((!mergeNbr1Allowed) && (!mergeNbr2Allowed)) {
    continue;
}

if (mergeNbr1Allowed &&
    ((nbr1->flags & NO_MESH_COARSEN) == 0)) {
    /* Existing MergeNode() attempt for nbr1. */
}

if (mergeDone == 0) {
    if (mergeNbr2Allowed &&
        ((nbr2->flags & NO_MESH_COARSEN) == 0)) {
        /* Existing MergeNode() attempt for nbr2. */
    }
}
\end{lstlisting}

\paragraph{Meaning.}
An edge touching a corner may be merged only if its minimum-image length is
$L\leq2r_{\mathrm{ann}}$.  The candidate guard skips a degree-two coarsen
candidate when both relevant corner edges are long, but retains an escape for
pathologically short segments.  At execution, each neighbor is gated
independently.  This prevents arbitrary arm ordering from collapsing a long
protected edge merely because the other edge is short.

\subsection{\texttt{src/RemoteOps.cc}}

\paragraph{Code added.}

\begin{lstlisting}[style=cpp]
void AddOpResetConstraint(Home_t *home, Tag_t *tag1, int constraint)
{
#ifdef _OPENMP
#pragma omp critical (CRIT_ADD_OP)
#endif
    {
        int                    opDataLen = sizeof(RemOpResetConstraint_t);
        int                    neededLen = home->opListUsedLen + opDataLen;
        RemOpResetConstraint_t opData;

        if (home->opListAllocLen < neededLen) {
            ExtendOpList(home);
        }

        opData.opType = REM_OP_RESET_CONSTRAINT;
        opData.tag1.domainID = tag1->domainID;
        opData.tag1.index    = tag1->index;
        opData.constraint    = constraint;

        memcpy(&home->opListBuf[home->opListUsedLen], &opData, opDataLen);
        home->opListUsedLen += opDataLen;
    }

    return;
}
\end{lstlisting}

\paragraph{Meaning.}
The emitter grows the operation buffer when necessary, packs the complete
constraint value, and appends the fixed-size record.  The same OpenMP critical
section used by other operation emitters prevents concurrent writers from
overlapping records.

\subsection{\texttt{src/FixRemesh.cc}}

\paragraph{Code added.}

\begin{lstlisting}[style=cpp]
} else if (opType == REM_OP_RESET_CONSTRAINT) {
    int                    opLen = sizeof(RemOpResetConstraint_t);
    RemOpResetConstraint_t opBuf;

    memcpy(&opBuf, &recvOpListBuf[offSet], opLen);
    offSet += opLen;

    node1 = GetNodeFromTag(home, opBuf.tag1);

    if (node1 != (Node_t *)NULL) {
        SET_CONSTRAINTS(node1->constraint, opBuf.constraint);
    }
\end{lstlisting}

\paragraph{Meaning.}
The receiver consumes exactly one matching wire record, resolves the target
tag, and replaces the complete constraint mask.  Ordering is essential:
\code{SplitNode()} emits create/connectivity operations first, then the caller
appends the reset.  Replay therefore creates the remote ghost before setting
\code{CORNER_NODE}, and it does so before immediate remeshing.

\section{Regression fixture}

\subsection{\texttt{tests/bcc\_binary\_junction\_node.ctrl}}

The control fixture is shown in full because it is both executable input and
the run guide.

\paragraph{File added.}

\begin{lstlisting}[style=data]
#
# One-step regression fixture for splitting a non-planar BCC binary-junction
# node.  Run from the tests directory with:
#
#   ../bin/paradis -d bcc_binary_junction_node.data \
#       bcc_binary_junction_node.ctrl
#
# The final restart must contain five nodes.  The original three-arm node is
# split into a three-arm node and a degree-two node carrying constraint 16
# (CORNER_NODE).  Set useParallelSplitMultiNode to 1 to exercise the parallel
# split implementation; 0 exercises the serial implementation.
# A two-rank communication check can then be run with:
#
#   mpirun -n 2 ../bin/paradis -doms 2 1 1 \
#       -d bcc_binary_junction_node.data bcc_binary_junction_node.ctrl
#
dirname = "bcc_binary_junction_node_results"

numXdoms = 1
numYdoms = 1
numZdoms = 1

numXcells = 4
numYcells = 4
numZcells = 4

fmEnabled       = 1
fmMPOrder       = 2
fmTaylorOrder   = 5
fmCorrectionTbl = "../inputs/fm-ctab.Ta.600K.0GPa.m2.t5.dat"

maxstep = 1
remeshRule = 2
minSeg = 2.000000e+02
maxSeg = 2.500000e+03
rann = 2.000000e+00
rTol = 1.000000e+00
timestepIntegrator = "trapezoid"

rc = 5.000000e+00
MobScrew = 1.000000e+01
MobEdge = 1.000000e+01
MobClimb = 1.000000e-08
mobilityLaw = "BCC_0"

split3node = 1
splitMultiNodeFreq = 1
useParallelSplitMultiNode = 0

loadType = 0
edotdir = [
  0.0
  0.0
  1.0
]
appliedStress = [
  0.000000e+00
  0.000000e+00
  5.000000e+08
  0.000000e+00
  0.000000e+00
  0.000000e+00
]

savetimers = 0
savecn = 1
savecnfreq = 1
savecncounter = 0
gnuplot = 0
povray = 0
saveprop = 0
\end{lstlisting}

\paragraph{Meaning.}
The fixture runs exactly one BCC step with a $zz$ stress of $5\times10^8$,
evaluates multi-node splitting every cycle, and writes a restart.  It defaults
to \code{useParallelSplitMultiNode = 0}, the serial topology implementation.
Set the selector to one in a \code{PARALLEL} build to exercise the distributed
force path.  The two-rank command overrides the domain layout to $2\times1
\times1$, testing ownership and remote constraint replay.

\subsection{\texttt{tests/bcc\_binary\_junction\_node.data}}

\paragraph{File added.}

\begin{lstlisting}[style=data]
dataFileVersion = 5
numFileSegments = 1
minCoordinates = [
  -1.750000e+04
  -1.750000e+04
  -1.750000e+04
]
maxCoordinates = [
   1.750000e+04
   1.750000e+04
   1.750000e+04
]
nodeCount = 4
dataDecompType = 2
dataDecompGeometry = [
  1
  1
  1
]

domainDecomposition =
# Dom_ID  Minimum XYZ bounds                    Maximum XYZ bounds
  0      -17500.0 -17500.0 -17500.0            17500.0 17500.0 17500.0

nodalData =
# A non-planar BCC binary junction: the <100> junction arm balances the two
# 1/2<111> glide arms.  The three fixed end nodes isolate the topology
# operation from endpoint motion.
 0,0 0.0 0.0 0.0 3 0
   0,1 -1.1547005383792517  0.0                 0.0
         0.0                 1.0                 0.0
   0,2  0.5773502691896258  0.5773502691896258  0.5773502691896258
        -0.7071067811865475  0.0                 0.7071067811865475
   0,3  0.5773502691896258 -0.5773502691896258 -0.5773502691896258
        -0.7071067811865475 -0.7071067811865475  0.0

 0,1 2000.0 0.0 0.0 1 7
   0,0  1.1547005383792517  0.0                 0.0
         0.0                 1.0                 0.0

 0,2 0.0 2000.0 0.0 1 7
   0,0 -0.5773502691896258 -0.5773502691896258 -0.5773502691896258
        -0.7071067811865475  0.0                 0.7071067811865475

 0,3 0.0 0.0 2000.0 1 7
   0,0 -0.5773502691896258  0.5773502691896258  0.5773502691896258
        -0.7071067811865475 -0.7071067811865475  0.0
\end{lstlisting}

\paragraph{Meaning.}
Node \code{0,0} is the unconstrained degree-three candidate at the origin.
The arm toward node \code{0,1} has a $\langle100\rangle$ Burgers-vector
magnitude of approximately 1.1547 and is recognized as the junction arm.  The
two unit-magnitude $\tfrac12\langle111\rangle$ arms balance it and form a
non-planar geometry under the five-degree test.  Nodes \code{0,1},
\code{0,2}, and \code{0,3} are 2000 units away and carry constraint 7
(fully pinned), isolating the topology event from endpoint motion.

\subsection{Commands and expected result}

Run from \path{tests}:

\begin{lstlisting}[style=data]
../bin/paradis -d bcc_binary_junction_node.data \
    bcc_binary_junction_node.ctrl
\end{lstlisting}

For the parallel topology selector, change
\code{useParallelSplitMultiNode = 1}.  For a two-rank ownership test, run:

\begin{lstlisting}[style=data]
mpirun -n 2 ../bin/paradis -doms 2 1 1 \
    -d bcc_binary_junction_node.data \
    bcc_binary_junction_node.ctrl
\end{lstlisting}

After one step,
\path{bcc_binary_junction_node_results/restart/restart.data} must contain five
nodes:

\begin{itemize}
  \item the three original pinned degree-one endpoints;
  \item exactly one degree-two node with constraint 16
        (\code{CORNER_NODE}); and
  \item one degree-three node without the corner constraint.
\end{itemize}

The regression review completed serial and \code{PARALLEL} builds, obtained
byte-identical one-rank restart data from selectors zero and one, completed
two-rank MPI runs with both selectors, and observed the expected unique
degree-two constraint-16 node.

\section{End-to-end meaning of the change}

The code fragments above form one ordered topology transaction:

\begin{enumerate}
  \item Parameter plumbing enables degree-three BCC candidates.
  \item The shared classifier identifies exactly one large-magnitude junction
        arm and rejects planar parent geometry.
  \item Serial or parallel arm enumeration moves only the two non-junction
        arms in the trial split.
  \item Mobility evaluation holds the degree-two trial node fixed, checks the
        separation direction, and accepts only a sufficiently positive power
        gain.
  \item The final split reuses the tested positions and marks the degree-two
        result as a corner.
  \item The append-only remote operation copies the full constraint state to
        ghost nodes before remeshing.
  \item Corner-aware coarsening preserves long physical edges while retaining
        a controlled escape for very short pathological segments.
\end{enumerate}

The change deliberately does not generalize all degree-three splitting.  It
is limited to BCC, exactly three valid arms, exactly one junction arm under the
Burgers-magnitude rule, non-planar geometry, and nodes that are neither pinned
nor surface-constrained.  Nodes with four or more arms retain their established
multi-node split behavior.

\end{document}
